\section{Evaluating \Veche{} on four pillars}
\label{sec:evaluation}

The Tournesol association proposed to evaluate digital democracy solutions
based on 4 digital democracy pillars.
According to these pillars, solutions must be:
\begin{description}
    \item[Transparent:] The solutions must be free and open-source,
    but also explainable and verifiable.
    \item[Secure:] The solutions must be resilient to coordinated attacks,
    especially through fake accounts. 
    When relevant, they should also allow for private weighing in
    (like voting in democracy).
    \item[Thoughtful:] The solutions must value informed and thoughtful judgments,
    as opposed to rushed, emotional and instinctive actions.
    \item[Uniting:] The solutions must unite people,
    by being inclusive and by contributing to a shared reality.
\end{description}
We refer to \url{tournesol.app/manifesto} for more justifications of the pillars.
We added a fifth evaluation based on user experience.
In this section, 
we evaluate the leading recommendation algorithms along the four pillars,
as well as \SimpleVeche{}, Tournesol-\Veche{} and BlueSky-\Veche{}.
The latter two include the generalizations discussed in Section~\ref{sec:generalizations},
with the hyperparameters provided in the various tables of the Section.
In particular, both leverage diversification
(which allows to be more uniting, 
as it values diverging viewpoints within the community of councillors).
Table~\ref{table:pillars} summarizes our analysis.
Its last line corresponds to the inclusion of common-ground recommendations,
as discussed in the following paragraph.

\begin{table}
    \centering
    \begin{tabular}{|c|c|c|c|c|c|}
        \hline
        & Transparent & Secure & Thoughtful & Uniting & UX \\ \hline
        Proprietary engagement maximization  & E & E & E & E & D \\ \hline
        Matrix factorization   & C & E & D & E & C \\ \hline
        Reverse chronological  & A & D & D & E & E \\ \hline
        \SimpleVeche{}         & A & C & D & E & D \\ \hline
        Tournesol-\Veche{}     & B & C & D & D & C \\ \hline
        BlueSky-\Veche{}       & B & C & D & D & C \\ \hline
        \Veche{} + Common*     & B & B & B & B & C \\ \hline
    \end{tabular}
    \caption{Evaluating Veche on Tournesol's four digital democracy pillars}
    \label{table:pillars}
\end{table}

\paragraph{Reinforcing unity.}
Given this analysis, we actually do not believe that \Veche{},
nor any follow-based algorithms,
should be viewed as a societally satisfactory solution.
In particular, such algorithms are inherently facilitating filter bubbles,
which may divide societies and harm democracies.
This will be even more the case if algorithmic representatives customize their ballots
based on the users they target.
To construct a common information ground that better unites citizens,
we believe that a fraction of the recommendations, say 30\%,
should instead result from a social choice mechanism that involves all citizens
(not only followed accounts),
and that yields non-personalized recommendations.
Our algorithms may in fact be instrumental
to both provide to the remaining 70\% of personalized recommendations
and to inspire mechanisms that respond to the need for common ground.
